\section*{Anhang}

Auf den Bildern sind teilweise die Pythonvariablen anstatt der Variablennamen zu erkennen. Eine entsprechende Übersetzung ist in Tabelle \ref{tab:feature_names} zu finden.

\begin{table}[H]
	\centering
    \begin{tabular}{|c|c|}
        \hline
        \textbf{Python-Variable} & \textbf{Feature-Name} \\
        \hline
        inst\_amp & Instantaneous Amplitude \\
        avg\_energy & Average Energy \\
        rms\_amp & RMS Amplitude \\
        quadrature & Quadrature\\
        inst\_q & Instantaneous Quality Factor \\
        inst\_phase\_real & Instantaneous Phase (Real Part) \\
        inst\_phase\_imag & Instantaneous Phase (Imaginary Part) \\
        inst\_freq & Instantaneous Frequency \\
        semblance & Semblance \\
        skewness & Skewness \\
        kurtosis & Kurtosis \\
        dip & Dip \\
        \hline
    \end{tabular}
    \caption{Übersicht der Python-Variablen und den zugehörigen Feature-Namen}
    \label{tab:feature_names}
\end{table}

\begin{table}[H]
    \centering
    \begin{tabular}{|c|c|c|c|c|c|c|}
        \hline
        \textbf{Merkmal} & $\mathbf{f_{min}}$ & $\mathbf{f_{max}}$ & $\mathbf{x_{mean}}$ & $\mathbf{t_{mean}}$ & $\mathbf{x_{window}}$ & $\mathbf{t_{window}}$ \\
        \hline
        inst\_amp     & 0     & 100   & 15        & 15        & -        & -        \\
        inst\_phase   & -     & -     & -         & -         & -        & -        \\
        inst\_freq    & -3e8  & 3e8   & 10 (100)  & 10 (100)  & -        & -        \\
        quadrature    & 0     & 50    & 40        & 10        & -        & -        \\
        inst\_q       & 0     & 10    & 10        & 10        & -        & -        \\
        avg\_energy   & 0     & 4e3   & 15        & 15        & 1        & 1        \\
        rms\_amp      & 0     & 300   & 15        & 15        & 1        & 1        \\
        semblance     & -     & -     & -         & -         & 40       & 20       \\
        skewness      & 0     & 10    & 60        & 40        & 1        & 1        \\
        kurtosis      & 0     & 5     & 60        & 40        & 1        & 1        \\
        dip           & -     & -     & 40        & 40        & -        & -        \\
        \hline
    \end{tabular}
    \caption{Parameterübersicht der Merkmale, Prozessierstufe 1\\
    $\mathbf{f_{min}}$ und $\mathbf{f_{max}}$ sind die Minima und Maxima des Wertebereichs. Wenn um ein Sample innerhalb eines Fensters gemittelt wurde, stehen $\mathbf{x_{mean}}$ und $\mathbf{t_{mean}}$ für die Fensterbstände zum Sample (Bsp.: $\mathbf{x_{mean}} = \mathbf{t_{mean}} = 3 \rightarrow$ 7 x 7 Fenster). Wird ein Merkmal auf einem Fenster berechnet, stehen $\mathbf{x_{window}}$ und $\mathbf{t_{window}}$ für die Fensterabstände zum Sample in x- und t-Richtung (analog wie $\mathbf{x_{mean}}$ und $\mathbf{t_{mean}}$). Die Werte in Klammern bei $\mathbf{x_{mean}}$ und $\mathbf{t_{mean}}$ beziehen sich auf verschiedene Modellläufe, bei denen weniger starkes und stärkeres Glätten genutzt wurde.}
    \label{tab:feature_params_1}
\end{table}

\begin{table}[H]
    \centering
    \begin{tabular}{|c|c|c|c|c|c|c|}
        \hline
        \textbf{Merkmal} & $\mathbf{f_{min}}$ & $\mathbf{f_{max}}$ & $\mathbf{x_{mean}}$ & $\mathbf{t_{mean}}$ & $\mathbf{x_{window}}$ & $\mathbf{t_{window}}$ \\
        \hline
        inst\_amp     & 0     & 300   & 15        & 15        & -        & -        \\
        inst\_phase   & -     & -     & -         & -         & -        & -        \\
        inst\_freq    & -3e8  & 3e8   & 10        & 10        & -        & -        \\
        quadrature    & -     & -     & -         & -         & -        & -        \\
        inst\_q       & 0     & 10    & 15        & 15        & -        & -        \\
        abs\_grad     & -     & -     & 15        & 15        & -        & -        \\
        avg\_energy   & 0     & 2e4   & 15        & 15        & 1        & 1        \\
        rms\_amp      & 0     & 300   & 15        & 15        & 1        & 1        \\
        semblance     & -     & -     & -         & -         & 40       & 20       \\
        kurtosis      & 0     & 5     & 60        & 40        & 1        & 1        \\
        dip           & -     & -     & 40        & 40        & -        & -        \\
        \hline
    \end{tabular}
    \caption{Parameterübersicht der Merkmale, Prozessierstufe 2\\
    $\mathbf{f_{min}}$ und $\mathbf{f_{max}}$ sind die Minima und Maxima des Wertebereichs. Wenn um ein Sample innerhalb eines Fensters gemittelt wurde, stehen $\mathbf{x_{mean}}$ und $\mathbf{t_{mean}}$ für die Fensterabstände zum Sample (Bsp.: $\mathbf{x_{mean}} = \mathbf{t_{mean}} = 3 \rightarrow$ 7 x 7 Fenster). Wird ein Merkmal auf einem Fenster berechnet, stehen $\mathbf{x_{window}}$ und $\mathbf{t_{window}}$ für die Fensterabstände zum Sample in x- und t-Richtung (analog wie $\mathbf{x_{mean}}$ und $\mathbf{t_{mean}}$).}
    \label{tab:feature_params_2}
\end{table}

\begin{table}[H]
    \centering
    \begin{tabular}{|c|c|c|c|c|c|c|}
        \hline
        \textbf{Merkmal} & $\mathbf{f_{min}}$ & $\mathbf{f_{max}}$ & $\mathbf{x_{mean}}$ & $\mathbf{t_{mean}}$ & $\mathbf{x_{window}}$ & $\mathbf{t_{window}}$ \\
        \hline
        inst\_amp     & 0     & 300   & 15        & 15        & -        & -        \\
        inst\_phase   & -     & -     & -         & -         & -        & -        \\
        inst\_freq    & -3e8  & 3e8   & 10        & 10        & -        & -        \\
        sweetness     & 0     & 2e-6  & 10        & 10        & -        & -        \\
        quadrature    & 0     & 50    & 40        & 10        & -        & -        \\
        inst\_q       & 0     & 10    & 15        & 15        & -        & -        \\
        abs\_grad     & -     & -     & 15        & 15        & -        & -        \\
        avg\_energy   & 0     & 4e3   & 15        & 15        & 1        & 1        \\
        rms\_amp      & 0     & 300   & 15        & 15        & 1        & 1        \\
        coherence     & -100  & 10    & 100       & 100       & 100      & 100      \\
        semblance     & -     & -     & -         & -         & 40       & 20       \\
        skewness      & 0     & 10    & 60        & 40        & 1        & 1        \\
        kurtosis      & 0     & 5     & 60        & 40        & 1        & 1        \\
        dip           & -     & -     & 40        & 40        & -        & -        \\
        \hline
    \end{tabular}
    \caption{Parameterübersicht der Merkmale, Prozessierstufe 3\\
    $\mathbf{f_{min}}$ und $\mathbf{f_{max}}$ sind die Minima und Maxima des Wertebereichs. Wenn um ein Sample innerhalb eines Fensters gemittelt wurde, stehen $\mathbf{x_{mean}}$ und $\mathbf{t_{mean}}$ für die Fensterabstände zum Sample (Bsp.: $\mathbf{x_{mean}} = \mathbf{t_{mean}} = 3 \rightarrow$ 7 x 7 Fenster). Wird ein Merkmal auf einem Fenster berechnet, stehen $\mathbf{x_{window}}$ und $\mathbf{t_{window}}$ für die Fensterabstände zum Sample in x- und t-Richtung (analog wie $\mathbf{x_{mean}}$ und $\mathbf{t_{mean}}$).}
    \label{tab:feature_params_3}
\end{table}

\begin{figure}[H]
    \centering
    \includegraphics[width=\textwidth]{feature_plots_raw/pairplot_overview.jpg}
    \caption{Pairplot der Rohdaten-Features, Prozessierstufe 1}
    \label{fig:pairplot_overview_raw}
\end{figure}

\begin{figure}[H]
    \centering
    \includegraphics[width=\textwidth]{feature_plots/pairplot_overview.jpg}
    \caption{Pairplot der Features, Prozessierstufe 2}
    \label{fig:pairplot_overview}
\end{figure}

\begin{figure}[H]
    \centering
    \includegraphics[width=\textwidth]{feature_plots_khigh/pairplot_overview.jpg}
    \caption{Pairplot der Features, Prozessierstufe 3}
    \label{fig:pairplot_overview_khigh}
\end{figure}